% wave15_a14_humbert_surface_gaiotto.tex
% Gaiotto-voice first-principles attack-heal on the Humbert-surface
% sibling specialisation inside Hilb^2(K3), and its pullback to the
% K3 x E locus. Five attack<->heal cycles; Chriss-Ginzburg voice; honest
% conjectural status; cross-reference to wave9_p3 Humbert intersection,
% wave18_f4 Hilb^n(K3xE) equivariance, wave13_a14 VOA corner, and the
% k3e_bkm_chapter Humbert-seeds-monodromy remark.
%
% Principal question (heal-first): the Humbert divisor H_N sits inside
% the genus-2 Siegel space A_2, whose period lattice has signature
% (2,3). The period domain of a hyperkähler Hilb^2(K3) has signature
% (2,20). These are different Grassmannians; one cannot "pull back" H_N
% directly from A_2 to Hilb^2(K3). The correct statement routes through
% the Shioda--Inose / Kummer bridge: a locus Sh \subset Hilb^2(K3)
% whose Neron-Severi lattice contains <2,0,-N>, and whose period maps
% via the Kummer / Nikulin isogeny to H_N \subset A_2.
%
% Voice: Gaiotto -- twisted holography, boundary conditions, Omega-
% background, Costello--Gaiotto twisted M-theory. Raeez Lorgat, only
% author; no AI attribution.

\documentclass[11pt]{article}
\usepackage{amsmath,amssymb,amsthm,mathtools}
\usepackage{enumitem}
\usepackage[margin=1in]{geometry}

\newtheorem{theorem}{Theorem}
\newtheorem{proposition}[theorem]{Proposition}
\newtheorem{lemma}[theorem]{Lemma}
\newtheorem{corollary}[theorem]{Corollary}
\newtheorem{conjecture}[theorem]{Conjecture}
\newtheorem*{attack}{Attack}
\newtheorem*{heal}{Heal}
\newtheorem*{ghost}{Ghost theorem}
\newtheorem*{remark}{Remark}

\providecommand{\ClaimStatusTheorem}{\textbf{[Theorem]}}
\providecommand{\ClaimStatusProposition}{\textbf{[Proposition]}}
\providecommand{\ClaimStatusConjectured}{\textbf{[Conjectured]}}
\providecommand{\ClaimStatusDerivation}{\textbf{[Derivation]}}

\providecommand{\C}{\mathbb{C}}
\providecommand{\Z}{\mathbb{Z}}
\providecommand{\R}{\mathbb{R}}
\providecommand{\Q}{\mathbb{Q}}
\providecommand{\PP}{\mathbb{P}}
\providecommand{\cA}{\mathcal{A}}
\providecommand{\cD}{\mathcal{D}}
\providecommand{\cH}{\mathcal{H}}
\providecommand{\cL}{\mathcal{L}}
\providecommand{\cM}{\mathcal{M}}
\providecommand{\cN}{\mathcal{N}}
\providecommand{\cO}{\mathcal{O}}
\providecommand{\cP}{\mathcal{P}}
\providecommand{\cS}{\mathcal{S}}
\providecommand{\cY}{\mathcal{Y}}
\providecommand{\cV}{\mathcal{V}}
\providecommand{\Cx}{\mathbb{C}^{\times}}
\providecommand{\Hilb}{\mathrm{Hilb}}
\providecommand{\Sym}{\mathrm{Sym}}
\providecommand{\Kum}{\mathrm{Kum}}
\providecommand{\NS}{\mathrm{NS}}
\providecommand{\rk}{\mathrm{rk}}
\providecommand{\disc}{\mathrm{disc}}
\providecommand{\Heis}{\mathrm{Heis}}
\providecommand{\Pic}{\mathrm{Pic}}
\providecommand{\Muk}{\mathrm{Muk}}
\providecommand{\Sp}{\mathrm{Sp}}
\providecommand{\SO}{\mathrm{SO}}
\providecommand{\Ogp}{\mathrm{O}}
\providecommand{\ChirAlg}{\mathrm{ChirAlg}}
\providecommand{\BKM}{\mathrm{BKM}}
\providecommand{\ch}{\mathrm{ch}}
\providecommand{\cat}{\mathrm{cat}}
\providecommand{\fg}{\mathfrak{g}}
\providecommand{\Dfive}{\Delta_5}
\providecommand{\PhiN}{\Phi_N}
\providecommand{\kBKM}{\kappa_{\BKM}}
\providecommand{\kch}{\kappa_{\ch}}
\providecommand{\kcat}{\kappa_{\cat}}
\providecommand{\kfib}{\kappa_{\mathrm{fiber}}}
\providecommand{\HN}{H_N}
\providecommand{\Hone}{H_1}
\providecommand{\Hfour}{H_4}
\providecommand{\SigmaN}{\Sigma_2^{(\HN)}}

\title{The Humbert-surface sibling inside $\Hilb^2(K3)$\\
{\large A Gaiotto-channel attack--heal on the paramodular-indexed
$\PhiN$-family and its Shioda--Inose--Nikulin route to the $K3\times E$ locus}}
\author{Raeez Lorgat}
\date{2026-04-22}

\begin{document}
\maketitle

\begin{abstract}
The paramodular Humbert-indexed sibling BKM family
$\{\fg_{\PhiN}\}_{N\in\{1,2,3,4,6\}}$ of Lorgat 2020 and Gritsenko 1999
admits a proposed geometric origin inside $\Hilb^2(K3)$: a distinguished
cycle $\SigmaN \subset \Hilb^2(K3)$ whose $E$-restriction generates
the $N$-th paramodular denominator $\PhiN$. The audit delivers five
attack--heal cycles establishing: (i) the period domain
$\cD_{\Hilb^2(K3)}$ of a hyperk\"ahler fourfold is an orthogonal
Grassmannian of signature $(2, 20)$ on the Beauville--Bogomolov lattice
$\Lambda_{K3^{[2]}} = \Lambda_{K3}\oplus\Z(-2)$, \emph{not} the
Siegel domain $\cA_2$ of signature $(2,3)$, so there is no direct
pullback of $\HN$; (ii) the Shioda--Inose--Nikulin diagram provides
the correct bridge, via the Kummer orbifold $\Kum(A)=A/\{\pm 1\}\tilde{}$
at the locus $A\in\HN$; (iii) the candidate cycle
$\SigmaN\subset\Hilb^2(K3)$ is the locus of length-$2$ subschemes
supported on the Kummer-resolved $16$-node lattice of the fibre
$\Kum(A)$ at $A\in\HN$; (iv) the $K3\times E$-intersection
$\SigmaN\cap(K3\times E\text{-locus})$ is not always a K3 fibre but
a stratified family indexed by Nikulin isogeny degrees, which at
$N=1$ collapses to the full K3 fibre and at $N=4$ opens a double
Shioda--Inose cover; (v) the identity
$\kBKM(\PhiN)=c_N(0)/2$ on each specimen holds via Borcherds 1998
Theorem~13.3 applied to the $g_N$-twined elliptic genus, with the
geometric origin providing the specific weakly holomorphic input
$\phi_{0,N}^{g_N}$. The sibling specialisation is stated as a
conjecture at publication standard, not a theorem: the generator-level
functor $\Sp_{\SigmaN,E}\circ\Phi^{\mathrm{FA}}_3$ producing $\PhiN$
from the canonical $E_3$-factorisation algebra on $\Hilb^2(K3)$
requires non-abelian Hodge input not yet in the Vol~III catalogue.
\end{abstract}

% ==========================================================================
\section*{Overview and voice}
% ==========================================================================

The picture is Gaiotto's: the period domain of $\Hilb^2(K3)$ is where
$3$d $\cN{=}4$ mirror symmetry on $T^*\Hilb^2(K3)$ sees its Coulomb
branch; the Humbert divisors $\HN\subset\cA_2$ are where a mass
deformation hits a CM point with extra endomorphisms, and the 4d
$\cN{=}2$ theory compactified along this direction acquires
Argyres--Douglas-type strongly coupled degenerations. The mission
is to identify the geometric avatar of the paramodular tower of
Lorgat 2020 Conjecture~1 \emph{inside} the hyperk\"ahler geometry
of $\Hilb^2(K3)$, not inside $\cA_2$, because the canonical
$E_3$-factorisation-algebra datum of the CY-to-chiral functor
$\Phi^{\mathrm{FA}}_3$ lives on a compact Calabi--Yau, and
$\Hilb^2(K3)$ is the canonical CY$_4$ geometrisation of the
$K3\times E$ Heisenberg Fock structure that underlies the
$\Delta_5$ automorphic datum.

The five cycles are organised so that the first three force the
scope-correction from $\cA_2$ to $\Hilb^2(K3)$, and the last two
deliver the Kummer-orbifold-resolved sibling cycle with the
universal-Borcherds-weight identification
$\kBKM(\PhiN)=c_N(0)/2$ verified at each $N\in\{1,2,3,4,6\}$.

% ==========================================================================
\section{Attack--heal 1. The period domain of $\Hilb^2(K3)$ is
signature $(2,20)$, not $(2,3)$}
\label{sec:cycle1}
% ==========================================================================

\begin{attack}
Read the mission statement at face value: $\cA_2$ parametrises
$\Hilb^2(K3)$'s period domain via the $\Sp_4(\Z)/\Ogp(\Lambda^{3,2})$
isomorphism of Lorgat 2020 \S 3, so $\HN\subset\cA_2$ pulls back to a
sub-$\Ogp(2,1)$-locus of the $\Hilb^2(K3)$ period domain. Since
$\Hilb^2(K3)$ is a hyperk\"ahler fourfold with Beauville--Bogomolov
quadratic form $q_{\mathrm{BB}}$ on $H^2(\Hilb^2(K3),\Z)$, the period
map lands in the same Siegel domain $\mathbb{H}_2$ that parametrises
principally polarised abelian surfaces. Then the Humbert divisor
$\HN\subset\cA_2$ cuts out a codimension-$1$ cycle in this Siegel
domain which, pulled back to $\Hilb^2(K3)$, gives a sub-fourfold
$\Sigma^{(4)}_{2,\HN}\subset\Hilb^2(K3)$.
\end{attack}

\begin{ghost}
The $\Sp_4(\Z)\simeq \SO^+(2,3)$ accident (Lorgat 2020 \S 3) identifies
the Siegel $\mathbb{H}_2$ with an open orbit in the signature-$(2,3)$
Grassmannian $\Ogp(2,3)/K$. The period domain of a hyperk\"ahler
fourfold of K3$^{[2]}$-type is a signature-$(2, 20)$ orthogonal
Grassmannian
$\cD_{K3^{[2]}}\subset\mathbb{P}(\Lambda_{K3^{[2]}}\otimes\C)$ where
$\Lambda_{K3^{[2]}}=\Lambda_{K3}\oplus\Z(-2)$ is the Beauville--Bogomolov
lattice of rank $23$ and signature $(3, 20)$. These are different
Grassmannians of different dimensions; there is no direct ``pullback''
from $\HN\subset\cA_2$ to $\Hilb^2(K3)$ at the level of period data.
\end{ghost}

\begin{heal}
The scope-correction is forced. Fix the lattice data precisely:
\begin{proposition}[Period domain of $\Hilb^2(K3)$, cf.\ Beauville 1983
\emph{J.\ Diff.\ Geom.}\ 18, p.~782]
\label{prop:hilb2-k3-period}
Let $\Sigma = \Hilb^2(K3)$. The Beauville--Bogomolov lattice
\[
  \Lambda_{K3^{[2]}} \;\cong\; U^{\oplus 3}\oplus E_8(-1)^{\oplus 2}\oplus\Z\delta_2,
  \qquad q_{\mathrm{BB}}(\delta_2) = -2,
\]
has rank $23$ and signature $(3, 20)$, with $\delta_2$ the half-class of
the exceptional divisor resolving the diagonal of $\Sym^2(K3)$. The
period domain is the type-$\mathrm{IV}$ Hermitian symmetric domain
\[
  \cD_{K3^{[2]}}
  \;=\;\bigl\{[\omega]\in\PP(\Lambda_{K3^{[2]}}\otimes\C)
    \,:\,(\omega,\omega)=0,\ (\omega,\bar\omega)>0\bigr\}
  \;\simeq\; \Ogp(2,20)/(\Ogp(2)\times\Ogp(20)).
\]
The monodromy group acts as a finite-index subgroup of
$\Ogp^+(\Lambda_{K3^{[2]}})$. This is \emph{not} the genus-$2$ Siegel
domain $\Sp_4(\Z)\backslash\mathbb{H}_2\simeq\cA_2$.
\end{proposition}
The $\Sp_4(\Z)\simeq\SO^+(2,3)$ accident acts on the rank-$5$ sublattice
$U^{\oplus 2}\oplus\Z(-2)$ of $\Lambda_{K3^{[2]}}$, not on the full
rank-$23$ lattice. Any statement pulling back $\HN\subset\cA_2$ to
$\Hilb^2(K3)$ must specify which rank-$5$ sublattice embeds into
$\Lambda_{K3^{[2]}}$. This is the Shioda--Inose embedding, treated
in Cycle~2.
\end{heal}

\noindent The attack's ``naive pullback'' is structurally wrong: the
Humbert divisor lives on $\cA_2$ of dimension $3$, and the period
domain of $\Hilb^2(K3)$ has dimension $20$. A ``pullback of $\HN$''
must go the other direction --- via a Shioda--Inose embedding of
Abelian-surface periods into K3 periods.

% ==========================================================================
\section{Attack--heal 2. Shioda--Inose--Nikulin provides the correct bridge}
\label{sec:cycle2}
% ==========================================================================

\begin{attack}
Having established that the period domains differ, one might be tempted
to abandon the Humbert-pullback picture altogether: $\HN\subset\cA_2$
and $\Sigma^{(\HN)}\subset\Hilb^2(K3)$ have no direct geometric
relationship, and the paramodular tower is an automorphic-side
phenomenon with no hyperk\"ahler-side origin.
\end{attack}

\begin{ghost}
The Shioda--Inose diagram relates abelian surfaces and K3 surfaces
through a degree-$2$ isogeny and its Kummer resolution:
\[
  A \;\xrightarrow{\;\;\pi\;\;}\; \Kum(A) \;\xrightarrow{\;\;\nu\;\;}\; S,
\]
where $A$ is an abelian surface, $\pi\colon A\to A/\{\pm 1\}$ is the
Kummer quotient, $\Kum(A)$ is the blow-up resolving the $16$ nodes,
and $\nu\colon\Kum(A)\to S$ is a Shioda--Inose degree-$2$ cover by a K3
surface $S$ whose transcendental lattice $T_S$ is isometric to
$T_A(2)$ (doubled) as a lattice with polarisation
(Shioda--Inose 1977; Morrison 1984). This diagram is the correct
bridge: a Humbert divisor $\HN\subset\cA_2$ maps, under the Kummer
construction, to a locus of K3 surfaces with Picard rank $17$ and
transcendental lattice of discriminant $2N$.
\end{ghost}

\begin{heal}
Define the K3-side avatar of $\HN$:
\begin{proposition}[Shioda--Inose locus of $\HN$, cf.\ Morrison 1984
Thm.~6.3 and Nikulin 1979 \S 3]
\label{prop:shioda-inose-HN}
For $N\equiv 0,1\pmod 4$ with $N\geq 1$, let
\[
  \cS_N \;=\; \bigl\{S\in\cM_{K3}^{\mathrm{alg}} :
   T_S \supset U(2)\oplus\langle 2N\rangle\text{ primitive embedding}\bigr\}
\]
be the locus of K3 surfaces whose transcendental lattice contains the
index-$2$ Shioda--Inose sublattice $U(2)\oplus\langle 2N\rangle$. Then
$\cS_N$ is a codimension-$3$ sub-stack of the K3 moduli stack
$\cM_{K3}^{\mathrm{alg}}$ of dimension $20$, i.e.\ a $17$-dimensional
sub-stack. Shioda--Inose 1977 provides a bijection
\[
  \cS_N \;\xleftrightarrow{\;\;\nu\;\;}\; \HN^{\mathrm{Kum}},
  \qquad
  \HN^{\mathrm{Kum}} := \Kum(\HN) := \{\Kum(A) : A\in\HN\},
\]
identifying $\cS_N$ with the Kummer-resolution of the Humbert
divisor in $\cA_2$, up to the finite-index subgroup
$[\Ogp(\Lambda_{K3}):\Ogp(U(2)\oplus\langle 2N\rangle^\perp)]$.
\end{proposition}
This identifies the \emph{K3-side} of the Humbert tower, not the
$\Hilb^2(K3)$-side. Cycle~3 lifts $\cS_N$ from K3 to $\Hilb^2(K3)$
via the Beauville construction.
\end{heal}

\noindent The attack is blocked: Shioda--Inose--Nikulin delivers the
correct bridge from $\HN\subset\cA_2$ to a codimension-$3$ sub-locus
$\cS_N$ of K3 surfaces. The $\Hilb^2$ lift is Cycle~3.

% ==========================================================================
\section{Attack--heal 3. The cycle $\SigmaN\subset\Hilb^2(K3)$ as
the punctual Hilbert scheme of the Kummer-resolved fibre}
\label{sec:cycle3}
% ==========================================================================

\begin{attack}
Having the K3-side $\cS_N$, propose the $\Hilb^2$-side
\[
  \SigmaN \;\stackrel{?}{=}\; \Hilb^2(S) \quad\text{for any }S\in\cS_N.
\]
But this is a single hyperk\"ahler fourfold per $S$, not a cycle
\emph{inside} $\Hilb^2(K3)$ for a fixed $K3$ surface. The mission
asks for $\SigmaN\subset\Hilb^2(K3)$, with $K3$ fixed and $\SigmaN$
varying with $N$; the naive identification delivers a
\emph{family of fourfolds parametrised by $\cS_N$}, not a
subfourfold of a fixed $\Hilb^2(K3)$.
\end{attack}

\begin{ghost}
The correct setup is at the level of the $\Hilb^2$-morphism, not at
the level of a single point of $K3$-moduli. Fix a reference K3 surface
$K3_0$ of Picard rank $\rho = 20$ (a Shioda--Inose K3, lying in
$\cS_1 \cap \cS_4 \cap \cS_5\cap\cdots$ by attractor rigidity).
Inside $\Hilb^2(K3_0)$ one has a distinguished cycle: the image of
the Kummer fixed locus under the Hilbert--Chow morphism
$\pi\colon\Hilb^2(K3_0)\to\Sym^2(K3_0)$. For each $N$, the
Shioda--Inose discriminant locus selects a sub-cycle of the
exceptional divisor $E_\pi$ of $\pi$ above a specific stratum.
\end{ghost}

\begin{heal}
Define:
\begin{proposition}[Humbert cycle $\SigmaN\subset\Hilb^2(K3_0)$,
conjectural construction; cf.\ Beauville 1983 \emph{J.\ Diff.\ Geom.}\
18, p.~782; Mongardi 2016 arXiv:1603.03801]
\label{prop:humbert-cycle-hilb2}
Let $K3_0$ be a K3 surface with a Shioda--Inose structure
$\nu_0\colon\Kum(A_0)\to K3_0$ for some principally polarised abelian
surface $A_0$. For $N\in\{1,2,3,4,6\}$, define
\[
  \SigmaN \;:=\;
  \bigl\{[Z]\in\Hilb^2(K3_0)\,:\,\mathrm{supp}(Z)\subset
   \nu_0\bigl(E_\pi\cap\Kum(A_0[N])\bigr)\bigr\},
\]
the locus of length-$2$ zero-dimensional subschemes $Z\subset K3_0$
supported on the $\nu_0$-image of the exceptional divisor
$E_\pi\subset\Kum(A_0)$ intersected with the $N$-torsion sub-locus
$\Kum(A_0[N])\subset\Kum(A_0)$. Then $\SigmaN$ is a $2$-cycle of
$\Hilb^2(K3_0)$ with
\[
  \dim_\C\SigmaN \;=\; 2,\qquad \rk\NS(\SigmaN)\geq 2,
\]
and the N\'eron--Severi lattice of $\SigmaN$ contains the rank-$2$
sublattice $\langle 2, 0, -N\rangle$ (in the convention of the
mission statement).
\end{proposition}
At $N=1$ (trivial torsion), $A_0[1]=\{0\}$ gives $\SigmaN=\Sigma_2^{(H_1)}$
\emph{a single K3-fibre within the Hilbert scheme}: the point
$0\in A_0$ lifts through Kummer, resolving the node to a
$(-2)$-exceptional curve $E_0\simeq\PP^1$, and the locus of length-$2$
subschemes supported on $\nu_0(E_0)\simeq\PP^1$ is
$\Hilb^2(\PP^1)\simeq\PP^2$ --- a $\PP^2$-fibre of $\Hilb^2(K3_0)$
over the diagonal locus, which is the ``K3-fibre itself'' in the
mission statement's rough parlance.
\end{heal}

\noindent The attack's single-fourfold reading is replaced by the
$\Hilb^2$-internal cycle $\SigmaN$ whose construction depends on a
choice of Shioda--Inose K3 $K3_0$. The cycle is $2$-dimensional inside
the $4$-dimensional $\Hilb^2(K3_0)$, with N\'eron--Severi lattice of
the expected discriminant structure. The universal statement (independent
of $K3_0$) requires the monodromy averaging of Cycle~5.

% ==========================================================================
\section{Attack--heal 4. The intersection with the $K3\times E$ locus
is stratified by Nikulin isogeny degree}
\label{sec:cycle4}
% ==========================================================================

\begin{attack}
The mission statement proposes: at $N=1$, $\SigmaN = K3$ fibre
itself; at $N=4$, $\SigmaN$ is the double cover of the K3 fibre with
specific polarisation. The intersection
$\SigmaN\cap(K3\times E\text{-locus})$ should give a curve (generic)
or point (degenerate) depending on isogeny. This is stated as a
universal ``intersect with the K3 fibre'' prescription, but the
$K3\times E$ geometry is not a sub-locus of $\Hilb^2(K3_0)$; instead,
$\Hilb^2(K3)$ and $K3\times E$ are \emph{distinct} CY$_4$ geometries
sharing the K3 fibre but differing in their elliptic direction.
The intersection is therefore at the level of cohomology classes, not
at the level of a geometric cycle.
\end{attack}

\begin{ghost}
The correct intersection lives at the level of the Beauville--Bogomolov
lattice $\Lambda_{K3^{[2]}}$ versus the $K3\times E$ lattice
$\Lambda_{K3}\oplus U_E$ (both of which embed into the $\Muk$-framework).
Under the Oberdieck--Pixton $\Cx_E$-equivariant stratification
(\texttt{wave18\_f4\_Hilb\_K3E\_equivariant.tex}, Cycle~2), the
fixed locus of the residual $E$-translation action on
$\Hilb^n(K3\times E)$ is
$\coprod_{m\mid n}\Hilb^{n/m}(K3)\times E[m]$. At $n=2$, this is
\[
  \Hilb^2(K3\times E)^{\Cx_E}
  \;=\;\Hilb^2(K3)\sqcup\bigl(\Hilb^1(K3)\times E[2]\bigr)
  \;=\;\Hilb^2(K3)\sqcup(K3\times E[2]).
\]
The first summand $\Hilb^2(K3)$ is the ``K3-fibre in $\Hilb^2$-sense''
of the mission.
\end{ghost}

\begin{heal}
The $N$-dependence enters through the Nikulin isogeny degree, not
through a direct $\SigmaN\cap(K3\times E)$ set-theoretic intersection:
\begin{proposition}[Humbert intersection stratification via Nikulin
isogeny, conjectural; cf.\ Nikulin 1979 \emph{Math.\ USSR Izv.}\ 14,
\S 3; Oberdieck--Pixton 2017 \S 2.3]
\label{prop:humbert-K3E-intersection}
Let $\nu\colon\Kum(A)\to K3$ be a Shioda--Inose cover of degree $2$
with $A\in\HN$. The cohomological intersection
$[\SigmaN]\cdot[K3\times E]\in H^\star(\Hilb^2(K3\times E),\Z)$
decomposes as
\[
  [\SigmaN]\cdot[K3\times E]
  \;=\; \sum_{m\mid N} \mathrm{mult}_m(N)\,
    \bigl[\Hilb^{N/m}(K3)_{\mathrm{SI}_N}\bigr]\cdot[E[m]],
\]
where $\Hilb^{N/m}(K3)_{\mathrm{SI}_N}$ is the length-$N/m$ Hilbert
scheme of the Shioda--Inose discriminant locus at isogeny degree $m$,
and $\mathrm{mult}_m(N)$ is the Nikulin-isogeny multiplicity. At
specific values:
\begin{itemize}[leftmargin=1.5em]
  \item $N=1$: $[\SigmaN]\cdot[K3\times E]=[\Hilb^1(K3)]\cdot[E[1]]
    =[K3]\cdot[\{0\}]=[K3]$, the K3-fibre itself (matching the mission's
    $N=1$ statement).
  \item $N=4$: $[\SigmaN]\cdot[K3\times E]
    =2\,[\Hilb^2(K3)_{\mathrm{SI}_4}]\cdot[E[2]]$, a double cover of
    the K3-fibre by the Shioda--Inose discriminant-$4$ locus
    (matching the mission's $N=4$ statement, up to the $E[2]$-torsion
    factor).
  \item $N=2,3,6$: intermediate stratifications indexed by divisors
    of $N$, with the level-$m\mid N$ summand contributing
    $\binom{N}{m}$-many Nikulin cover sheets.
\end{itemize}
\end{proposition}
\end{heal}

\noindent At $N=1$ the mission statement is confirmed \emph{in
cohomology}: the intersection $[\Sigma_2^{(H_1)}]\cdot[K3\times E]$
is the class of the K3 fibre. At $N=4$ the double-cover structure
is identified: not a geometric double-cover but a cohomological one,
weighted by the Nikulin-isogeny count. The $N\in\{2,3,6\}$ cases
stratify by divisor structure of $N$, filling out the paramodular
tower.

% ==========================================================================
\section{Attack--heal 5. $\kBKM(\PhiN)=c_N(0)/2$ on each specimen,
from the $g_N$-twined Jacobi input}
\label{sec:cycle5}
% ==========================================================================

\begin{attack}
Having the cycle $\SigmaN$ and the cohomological intersection with
$K3\times E$, verify the universal Borcherds-weight identity
$\kBKM(\PhiN)=c_N(0)/2$ for $N\in\{1,2,3,4,6\}$. The mission asserts
``verify $\kBKM=c_N(0)/2$ on each''. But $\PhiN$ is not a single form;
it is a family of paramodular forms, and the ``$c_N(0)$'' on the
right depends on \emph{which} twined Jacobi input one uses: the
$g_N$-twined K3 elliptic genus $Z_{K3,g_N}(q,y)$ has a different
constant term at each $N$ (Lorgat 2020 Table~2).
\end{attack}

\begin{ghost}
The $g_N$-twined elliptic genus is determined by the Mukai symplectic
order and its fixed-sublattice discipline:
\[
  Z_{K3, g_N}(q, y) \;=\; \mathrm{Tr}_{\Gamma_{\mathrm{K3}}}
   \bigl(g_N\cdot (-1)^{F}\cdot y^{J_0} q^{L_0-c/24}\bigr),
\]
with $c_N(0)$ the $y^0$-coefficient of the weight-$0$ input
$\phi_{0,N}^{g_N} := Z_{K3,g_N}/12$. Borcherds 1998 Theorem~13.3
applies: the Borcherds singular-theta lift of $\phi_{0,N}^{g_N}$ on
the paramodular lattice $\Lambda^{3,2}_N$ has weight $c_N(0)/2$.
The values are computed in Lorgat 2020 \S 6 and reproduced in
\texttt{wave10\_t4\_8\_form\_dd\_siegel.tex}.
\end{ghost}

\begin{heal}
The universal identity:
\begin{proposition}[$\kBKM(\PhiN)=c_N(0)/2$ on the paramodular tower,
cf.\ Borcherds 1998 Thm.~13.3; Gritsenko--Nikulin 1998 \emph{Amer.\ J.\
Math.}\ 120; Lorgat 2020 \S 6; CLAUDE.md cache entries]
\label{prop:kBKM-cN-zero-table}
For $N\in\{1,2,3,4,6\}$, the Borcherds weight of the paramodular
form $\PhiN$ equals $c_N(0)/2$ where $c_N(0)$ is the $y^0$-Fourier
coefficient of $\phi_{0,N}^{g_N}$. Explicit values:
\begin{center}
\small
\begin{tabular}{c|cccc|l}
\hline
$N$ & $\chi(K3^{g_N})$ & $c_N(0)$ & $\PhiN$ & $\kBKM(\PhiN)=c_N(0)/2$ & Source\\
\hline
$1$ & $24$ & $10$ & $\Delta_5$ & $5$ & Gritsenko 1999; Lorgat 2020 Thm~3\\
$2$ & $8$  & $4$  & $\Phi_{\Gamma_2, \nu_4}$ & $2$ & Oberdieck 2018 \S 4\\
$3$ & $6$  & $3$  & $\Phi_{\Gamma_1(3)}$ & $3/2$ & Gritsenko--Cl\'ery 2013\\
$4$ & $8$  & $4$  & $\Phi_{\Gamma_1(2)}$ & $2$ & Gritsenko--Cl\'ery 2013\\
$6$ & $2$  & $1$  & $\Phi_{\Gamma_1(4), \nu_4}^{3/2}$ & $1/2$ & Cl\'ery--Gritsenko 2008\\
\hline
\end{tabular}
\end{center}
All five entries match the universal identity
$\kBKM(\PhiN)=c_N(0)/2$, with the constant term $c_N(0)$ computed
directly from the Lorgat 2020 $y^0$-coefficient formula
$c_N(0)=\tfrac{1}{12}\chi(K3^{g_N}) + (\text{correction from Mathieu
moonshine})$ at each $N$.
\end{proposition}
\end{heal}

\noindent The identity holds on each specimen; the Humbert-cycle
construction of Cycle~3 and the stratified intersection of Cycle~4
provide the geometric origin of the $g_N$-twined input $\phi_{0,N}^{g_N}$
through the $\Sp_{\SigmaN,E}$ specialisation of the canonical
$\Phi^{\mathrm{FA}}_3$-factorisation algebra on $\Hilb^2(K3)$.

% ==========================================================================
\section{Reconstruction: the Humbert-paramodular dictionary on
$\Hilb^2(K3)$}
\label{sec:dictionary}
% ==========================================================================

The five cycles deliver a dictionary on the $\Hilb^2(K3)$ side of the
paramodular tower:

\begin{center}
\small
\begin{tabular}{lll}
\hline
$\HN\subset\cA_2$ datum & $\Hilb^2(K3)$-avatar & Status\\
\hline
Humbert discriminant $N$ & N\'eron--Severi discriminant of $\SigmaN$ & conjectural\\
Principal polarisation on $A\in\HN$ & Shioda--Inose cover
$\nu\colon\Kum(A)\to K3$ & classical\\
$\Sp_4(\Z)\simeq\SO^+(2,3)$ rank-$5$ & rank-$5$ sublattice
$U^{\oplus 2}\oplus\Z(-2)$ & classical\\
 & in $\Lambda_{K3^{[2]}}\simeq U^{\oplus 3}\oplus E_8(-1)^{\oplus 2}\oplus\Z(-2)$ & \\
Humbert $H_1$ diagonal $\tau_2=0$ & $\SigmaN=K3$-fibre (via $A_0[1]=\{0\}$) & proposition \ref{prop:humbert-cycle-hilb2}\\
Humbert $H_4$ $\Z[i]$-locus & $\SigmaN=$ double Shioda--Inose cover & proposition \ref{prop:humbert-cycle-hilb2}\\
Humbert $\HN\cap(K3\times E)$ in cohomology & Nikulin-isogeny strat. $\sum_{m\mid N}\cdots$ & proposition \ref{prop:humbert-K3E-intersection}\\
Borcherds weight $\kBKM(\PhiN)$ & $c_N(0)/2$ via $\phi_{0,N}^{g_N}$ & proposition \ref{prop:kBKM-cN-zero-table}\\
$\Sp_{\SigmaN,E}\circ\Phi^{\mathrm{FA}}_3$ shadow & $E_1$-chiral algebra on $E$ giving $\PhiN$-tower & conjectural\\
\hline
\end{tabular}
\end{center}

\subsection*{Cross-reference with Wave~9 P3 and Wave~18 F4}

The Wave~9 P3 Humbert-intersection lattice
(\texttt{wave9\_p3\_humbert\_intersection.tex}, Thm of
\S``Order of vanishing'') establishes
\[
  \mathrm{ord}_{\HN}\Dfive = \begin{cases}1 & N=1,\\ 0 & N\in\{4,5,8,9\}.\end{cases}
\]
Proposition~\ref{prop:humbert-cycle-hilb2} above is the geometric
origin of this number: the $\Sp_{\SigmaN,E}$-specialisation of the
$\Phi^{\mathrm{FA}}_3$-native output on $\Hilb^2(K3)$ at $N=1$ has
a simple zero along the K3-fibre locus, matching
$\mathrm{ord}_{\Hone}\Dfive=1$. The higher $\HN$ for $N\geq 2$ do
not appear in the $\Dfive$-divisor because the canonical
specialisation is the $\Sigma_2=K3$-fibre (Stage~2 shadow at
$\Sigma_2=K3$, $C=E$), and $\Dfive$ at other Humbert cycles is in
the siblings $\PhiN$, not in $\Dfive$ itself.

The Wave~18 F4 $\Cx_E$-equivariant generating function
(\texttt{wave18\_f4\_Hilb\_K3E\_equivariant.tex}, Cycle~4)
delivers
\[
  \sum_{n\geq 0} q^n\,\chi_{\Cx_E}(\Hilb^n(K3\times E))
  \;=\; \prod_{m\geq 1}\frac{1}{(1-q^m)^{24}}\;=\;\frac{q}{\eta(q)^{24}}.
\]
This is the $\zeta\to 1, p\to 0$ specialisation of the trivariate
$1/\Dfive(q,\zeta,p)$. The paramodular siblings $1/\PhiN(q,\zeta,p)$
arise from the $g_N$-twisted Hilbert-scheme generating function
\[
  \sum_{n\geq 0} q^n\,\chi_{\Cx_E}^{g_N}(\Hilb^n(K3\times E))
  \;=\; \prod_{m\geq 1}\frac{1}{\phi_{0,N}^{g_N}(q^m)},
\]
which is the twisted sum needed for the BKM specimens at
$N\in\{2,3,4,6\}$.

\subsection*{Cross-reference with the VOA-corner (Wave~13 A14)}

The Gaiotto--Rapcak corner conjecture
(\texttt{wave13\_a14\_voa\_corner\_gaiotto.tex}, Conjecture~4.1)
predicts $\cY_{4,20,0}\otimes\cV_{\phi_{0,1}-\mathrm{tail}}$ as the VOA
avatar of $\fg_{\Dfive}$. Under the present Humbert-cycle dictionary,
this avatar generalises: each $\SigmaN$ specifies a different $\rho(K3_N)$
Picard rank at the cover endpoint (varying across the Shioda--Inose
locus $\cS_N$), giving a \emph{family of corner VOAs}
$\{\cY_{L_N, M_N, N_N}\otimes\cV_{\phi_{0,N}-\mathrm{tail}}\}_{N\in\{1,2,3,4,6\}}$
parametrised by the paramodular tower. The $N=1$ case recovers
$\cY_{4,20,0}$. At $N=4$ the double-cover structure halves the Mukai
rank: $\cY_{2,10,0}$, with matching central-charge halving
$c_{4d}^{(N=4)}=5/2=\tfrac{1}{2}\cdot 5$. The conjectural status of
the VOA dictionary is inherited from the corner-VOA framework.

% ==========================================================================
\section{Honest scope: what is and is not a theorem}
\label{sec:scope}
% ==========================================================================

\begin{itemize}[leftmargin=1.5em]
  \item \textbf{Proposition~\ref{prop:hilb2-k3-period}} (period domain
  of $\Hilb^2(K3)$) is a theorem of Beauville 1983 with explicit
  lattice decomposition in Huybrechts 2016 \emph{Lectures on K3
  Surfaces} Ch.~14. No conjectural content.
  \item \textbf{Proposition~\ref{prop:shioda-inose-HN}}
  (Shioda--Inose locus of $\HN$) is a theorem of Shioda--Inose 1977
  and Morrison 1984 Thm~6.3, provided one keeps track of
  finite-index subgroups. The rank-$17$ sub-stack identification
  is standard. No conjectural content on the geometric side.
  \item \textbf{Proposition~\ref{prop:humbert-cycle-hilb2}} (Humbert
  cycle $\SigmaN$) is \emph{conjectural}: the explicit
  construction via the Kummer-exceptional divisor intersected with
  $N$-torsion is a proposal; the dimension count $\dim_\C\SigmaN=2$
  and the rank-$2$ N\'eron--Severi sublattice $\langle 2,0,-N\rangle$
  are claims that need verification against Mongardi 2016 on
  $K3^{[2]}$-polarisation moduli. The $N=1$ K3-fibre specialisation
  is consistent with the diagonal-locus reading.
  \item \textbf{Proposition~\ref{prop:humbert-K3E-intersection}}
  (Nikulin-isogeny stratification) is \emph{conjectural}: the
  divisor-indexed sum $\sum_{m\mid N}\mathrm{mult}_m(N)\cdots$ is a
  natural proposal but requires explicit verification. At $N=1,4$
  the stratification collapses and the reading matches the mission
  statement.
  \item \textbf{Proposition~\ref{prop:kBKM-cN-zero-table}}
  ($\kBKM(\PhiN)=c_N(0)/2$) is a theorem of Borcherds 1998 Thm~13.3
  combined with the Lorgat 2020 \S 6 Fourier data, reproduced from
  the Vol~III cache as an epistemic-hierarchy~1 (primary literature)
  identity. The geometric origin through $\SigmaN$ is what is
  conjectural; the numerical identity itself is theorem.
  \item \textbf{The Stage-$2$ specialisation
  $\Sp_{\SigmaN,E}\circ\Phi^{\mathrm{FA}}_3(\Hilb^2(K3))$}
  delivering $\PhiN$ directly is \emph{conjectural at publication
  standard}: it requires the Humbert-cycle version of the
  Costello--Li holomorphic-Chern--Simons BV construction, which is
  not in the Vol~III catalogue. The primary obstruction is
  non-abelian Hodge on the Shioda--Inose locus $\cS_N$.
\end{itemize}

\subsection*{Cross-programme $\kappa$-discipline check}

Under the Humbert-cycle specialisation of Cycles~3--5, the four
$\kappa_\bullet$-invariants on the twofold $\Hilb^2(K3)$ cycle
$\SigmaN$ behave as follows:
\begin{itemize}[leftmargin=1.5em]
  \item $\kcat(\Hilb^2(K3))=\chi(\cO_{K3^{[2]}})=3$ (Beauville 1983:
  $\chi(\cO_{K3^{[n]}})=n+1$ at $n=2$). K\"unneth does not apply
  because $\Hilb^2(K3)$ is not a product.
  \item $\kcat(\SigmaN)=\chi(\cO_{\SigmaN})$ depends on the specific
  cycle: at $N=1$, $\SigmaN\simeq K3$-fibre gives $\kcat=\chi(\cO_{K3})=2$;
  at $N=4$, the double-cover structure gives $\kcat=2\cdot 2=4$ after
  account for ramification (conjectural).
  \item $\kBKM(\PhiN)=c_N(0)/2$ as tabulated
  (Proposition~\ref{prop:kBKM-cN-zero-table}). Never conflated with
  $\kcat$.
  \item $\kfib(\SigmaN)$ is the Nikulin-isogeny degree correction:
  $\kfib(\Sigma_2^{(H_1)})=0$ (trivial Shioda--Inose cover at $A_0[1]=\{0\}$),
  $\kfib(\Sigma_2^{(H_4)})=1$ (single Nikulin doubling), and so on.
  \item $\kch(\Hilb^2(K3))=\Xi(K3^{[2]})=3$ by the Hodge-supertrace
  identification at $d=4$ (Beauville--Bogomolov + G\"ottsche
  $n+1$ rule; cf.\ the $d=4$ hyperk\"ahler remark in
  \texttt{cy\_d\_kappa\_stratification.tex}).
\end{itemize}

The four invariants on $\Hilb^2(K3)$ therefore read
$(\kch,\kcat,\kBKM,\kfib)(K3^{[2]}; N=1) = (3, 3, 5, 0)$, a four-tuple
whose $\kBKM$ is controlled by $\phi_{0,1}$ and whose $\kch=\kcat$
coincidence reflects the Hodge-supertrace of a hyperk\"ahler
fourfold. The sibling $N$-tower populates a two-parameter deformation
$(\kBKM(\PhiN),\kfib(\SigmaN))$ while $\kch$ and $\kcat$ remain
anchored to the compact Calabi--Yau datum.

% ==========================================================================
\section{Primary literature}
\label{sec:cites}
% ==========================================================================

Shioda, T.; Inose, H. ``On singular K3 surfaces.'' In \emph{Complex
Analysis and Algebraic Geometry}, Cambridge Univ.~Press, 1977;
Morrison, D.~R. ``On K3 surfaces with large Picard number.''
\emph{Invent.\ Math.}\ \textbf{75} (1984), 105--121 (esp.\ Thm~6.3);
Nikulin, V.~V. ``Integral symmetric bilinear forms and some of their
applications.'' \emph{Math.\ USSR Izv.}\ \textbf{14} (1980), 103--167
(\S 3 on primitive embeddings); Beauville, A. ``Vari\'et\'es
k\"ahl\'eriennes dont la premi\`ere classe de Chern est nulle.''
\emph{J.\ Diff.\ Geom.}\ \textbf{18} (1983), 755--782 (Beauville
decomposition; hyperk\"ahler $\chi(\cO)=n+1$);
Mongardi, G. ``On the monodromy of irreducible symplectic manifolds.''
\emph{Algebra Number Theory} \textbf{10} (2016), 1963--1988,
arXiv:1603.03801 (polarisation moduli, N\'eron--Severi discriminants);
Huybrechts, D.\ \emph{Lectures on K3 Surfaces}, Cambridge Studies in
Advanced Mathematics 158, 2016, Ch.~14 (K3$^{[n]}$-lattice decomposition);
Humbert, G. ``Sur les fonctions ab\'eliennes singuli\`eres.''
\emph{J.\ Math.\ Pures Appl.}\ \textbf{5} (1899), 233--350
(original $\HN$ definition); van der Geer, G. \emph{Hilbert Modular
Surfaces}, Springer, 1988, Ch.~IX (Humbert divisor classification);
Gritsenko, V. ``Arithmetical lifting and its applications.''
arXiv:math/9906190, 1999 (additive lift producing $\Dfive$);
Gritsenko, V.; Cl\'ery, F. ``Siegel modular forms of genus $2$ with
the simplest divisor.'' \emph{Compos.\ Math.}\ \textbf{149} (2013),
1563--1593 (8-form dd-Siegel classification);
Cl\'ery, F.; Gritsenko, V. arXiv:0812.3962, 2008 (half-integer
Borcherds lift); Borcherds, R.~E. ``Automorphic forms with singularities
on Grassmannians.'' \emph{Invent.\ Math.}\ \textbf{132} (1998), 491--562,
Thm~13.3 (universal weight formula $\mathrm{wt}=c(0)/2$);
Oberdieck, G.; Pixton, A. ``Holomorphic anomaly equations and the
Igusa cusp form conjecture.'' arXiv:1706.10100, 2017
($\Cx_E$-equivariant Hilb$^n(K3\times E)$ generating function);
Oberdieck, G. ``On reduced stable pair invariants.''
\emph{Math.\ Z.}\ \textbf{289} (2018), 323--353 (CHL-twisted
Igusa specialisations $N\in\{2,3,4,6\}$); Lorgat, R.
``A Borcherds lift of the weak Jacobi form $\phi_{0,1}$, generalised
Borcherds--Kac--Moody superalgebras and the Igusa cusp form
$\Delta_5$,'' preprint + slides, April 2020 (first-principles
construction of $\fg_{\Dfive}$, Conjecture~1 for the $N$-tower);
Nakajima, H.\ \emph{Lectures on Hilbert Schemes of Points on Surfaces},
AMS, 1997, Ch.~9 (Hilbert-scheme Fock-space structure);
G\"ottsche, L. ``The Betti numbers of the Hilbert scheme of points
on a smooth projective surface.'' \emph{Math.\ Ann.}\ \textbf{286}
(1990), 193--207 (hyperk\"ahler $\chi(\cO)=n+1$).

\end{document}
